\section*{Hoofdstuk \ref{ch:b3:gfd-foundations} --- Geofysische vloeistofdynamica}

\begin{solution}{exo:b3:gfd-foundations:1}
(a)~$f = 2\Omega\sin\varphi$: \num{7.292e-5}, \num{1.031e-4} en
\qty{1.263e-4}{s^{-1}}; $2\pi/f$: 23.93, 16.92 en \qty{13.82}{h}.
(b)~$U/f$: 2.74, 1.94 en \qty{1.58}{km}. (c)~Bij \ang{30}, $\sin\varphi =
1/2$ en $2\pi/f = 2\pi/\Omega$: de traagheidsperiode is de sterrendag,
\qty{23.93}{h}.
\end{solution}

\begin{solution}{exo:b3:gfd-foundations:2}
Met $f = \qty{1.031e-4}{s^{-1}}$: stofduivel \num{9.7e3}, zeebries 2.4,
depressie 0.15, gyre \num{2.4e-4}. De depressie en de gyre liggen dicht bij
geostrofisch evenwicht; de stofduivel draait beide kanten op, onbewust van de
rotatie van de Aarde, en de zeebries voelt het pas na een dag.
\end{solution}

\begin{solution}{exo:b3:gfd-foundations:3}
$\partial u_g/\partial z = (g/fT)\,|\partial T/\partial y| = 9.81\times
\num{5e-6}/(\num{1.031e-4}\times250) = \qty{1.90e-3}{s^{-1}}$; over
\qty{10}{km} neemt de westelijke wind met \qty{19}{m/s} toe.
\end{solution}

\begin{solution}{exo:b3:gfd-foundations:4}
$E = \num{50e-8}$: $N = \qty{2.21e-3}{s^{-1}}$, een periode van \qty{47}{min},
\num{1.27} cycli per uur. $E = \num{1000e-8}$: $N = \qty{9.90e-3}{s^{-1}}$,
\qty{10.6}{min}, 5.7 cycli per uur. $E = \num{e-8}$: $N =
\qty{3.13e-4}{s^{-1}}$, \qty{5.6}{h}, 0.18 cycli per uur.
\end{solution}

\begin{solution}{exo:b3:gfd-foundations:5}
(a)~$W \sim UH/L = \qty{0.1}{m/s}$ en $\mathrm Dw/\mathrm Dt \sim UW/L =
\qty{e-6}{m/s^2}$, \num{1.0e-7} van $g$. De drukterm, $fUL/H =
\qty{0.103}{m/s^2}$, is \num{1e5} keer groter: de verhouding, \num{9.7e-6}, is
$\mathrm{Ro}(H/L)^2$ met $\mathrm{Ro} = 0.097$. (b)~$\tilde fu =
2\Omega\cos\ang{45}\times10 = \qty{1.03e-3}{m/s^2}$, \num{1.05e-4} van $g$.
\end{solution}

\begin{solution}{exo:b3:gfd-foundations:6}
$fr = \qty{30}{m/s}$. (a)~$\mathrm{Ro}_g = 0.667$, $V = 40/(1 + \sqrt{3.667})
= \qty{13.7}{m/s}$, een derde langzamer dan geostrofisch. (b)~$\mathrm{Ro}_g =
0.167$, $V = 10/(1 + \sqrt{0.333}) = \qty{6.34}{m/s}$. (c)~$V_g \le fr/4 =
\qty{7.5}{m/s}$, en dan $V = fr/2 = \qty{15}{m/s}$, tweemaal de geostrofische
snelheid.
\end{solution}

\begin{solution}{exo:b3:gfd-foundations:7}
(a)~$\Omega = 2\pi/6 = \qty{1.047}{rad/s}$, $f = \qty{2.094}{s^{-1}}$.
$\mathrm{Ro} = 0.01/(2.094\times0.05) = 0.095$ en $\mathrm{Ek} =
\num{e-6}/(2.094\times0.04) = \num{1.19e-5}$: beide klein, dus de stroming voldoet aan
de stelling van Taylor--Proudman en er vormt zich een \omterm{thm:b3:gfd-foundations:taylor-proudman}{Taylorkolom}. (b)~$0.2/\sqrt{\num{e-6}
\times1.047} = \qty{195}{s}$, \qty{3.3}{min}, 33 omwentelingen van de draaitafel.
\end{solution}

\begin{solution}{exo:b3:gfd-foundations:8}
\sloppy $\partial\bar T/\partial y = -f\,\Delta u/(R\ln(850/250))$, met
$\ln(850/250) = 1.224$. Op \ang{30}~NB, $-\num{7.29e-5}\times36.12/(286.9
\times1.224) = \qty{-7.50}{K}$ per \qty{1000}{km}; op \ang{45}~NB,
$-\num{1.031e-4}\times17.18/351.1 = \qty{-5.05}{K}$ per \qty{1000}{km}. De
windschuiving is evenredig met de gradiënt gedeeld door $f$: de gradiënt is 1.5
maal groter op \ang{30} en $f$ 1.4 maal kleiner, dus de windschuiving is 2.1 maal
groter.\par
\end{solution}

\begin{solution}{exo:b3:gfd-foundations:9}
(a)~$L_R = 0.0113\times11\,000/\num{1.031e-4} = \qty{1206}{km}$;
$0.00357\times1000/\num{9.47e-5} = \qty{37.7}{km}$, 32 maal kleiner.
(b)~$\beta = 2\Omega\cos\varphi/a$: \num{1.62e-11} en
\qty{1.74e-11}{m^{-1}.s^{-1}}. $\beta L_R/f = 0.19$ voor de atmosfeer, 0.007
voor de oceaan.
\end{solution}

\begin{solution}{exo:b3:gfd-foundations:10}
(a)~$\mathrm dw/\mathrm dt = -(r + \mathrm if)w$, dus $w = U\mathrm
e^{-rt}\mathrm e^{-\mathrm ift}$: de snelheid draait met tempo $f$, terwijl
de grootte ervan vervalt als $\mathrm e^{-rt}$. De positie $x + \mathrm iy =
U(1 - \mathrm e^{-(r + \mathrm if)t})/(r + \mathrm if)$ spiraalt naar het
punt $U/(r + \mathrm if)$. (b)~$x + \mathrm iy = U(1 - \mathrm e^{-(r +
\mathrm if)t})/(r + \mathrm if) + u_0t$: de spiraal wordt uitgetrokken tot lussen.
(c)~$0.2\,\mathrm e^{-1} = \qty{0.074}{m/s}$; het centrum is $u_0t =
0.05\times4\times86\,400 = \qty{17.3}{km}$ oostwaarts verschoven.
\end{solution}

\begin{solution}{exo:b3:gfd-foundations:11}
$f = \qty{9.37e-5}{s^{-1}}$; $\partial u/\partial z = (g/f\rho_0)\,\partial
\rho/\partial y = 9.81\times10^{-5}/(\num{9.37e-5}\times1025) =
\qty{1.02e-3}{s^{-1}}$. Met $u = 0$ op \qty{1000}{m} is de oppervlaktestroom
\qty{1.02}{m/s}, oostwaarts: op het noordelijk halfrond stroomt de stroming
met het lichte water aan haar rechterkant.
\end{solution}

\begin{solution}{exo:b3:gfd-foundations:12}
(a)~$N^2 = (g/\rho)\,\Delta\sigma_\theta/\Delta z$. 0--\qty{100}{m}:
$9.81\times0.668/(1026.4\times100)$, $N = \qty{7.99e-3}{s^{-1}}$, 4.6
cycli per uur; 100--\qty{500}{m}: \qty{3.02e-3}{s^{-1}}, 1.7 cycli per
uur; 500--\qty{1000}{m}: \qty{3.24e-3}{s^{-1}}, 1.9 cycli per uur.
(b)~De bovenste \qty{100}{m}: deze bevat het warme oppervlaktewater boven de
seizoenspycnoklien. Tussen 100 en \qty{500}{m} is het water
gelijkmatiger dan eronder. (c)~$E = N^2/g$ is \num{650e-8}, \num{93e-8} en
\qty{107e-8}{m^{-1}}, allemaal binnen het bereik van de bovenste oceaan.
\end{solution}

\begin{solution}{pb:b3:gfd-foundations:1}
\textbf{1.} $f = 2\times\num{7.292115e-5}\sin\ang{48} =
\qty{1.084e-4}{s^{-1}}$; $2\pi/f = \qty{16.1}{h}$.

\textbf{2.} Badkuip \num{3.1e3}; tank ($f = 4\pi/10 = \qty{1.257}{s^{-1}}$)
0.020; bries 2.3; depressie 0.14; gyre \num{2.3e-4}.

\textbf{3.} De badkuip voelt niets van de rotatie van de Aarde; de bries
voelt het pas na enkele uren. De tank, de depressie en de gyre hebben
kleine Rossby-getallen en zijn dicht bij geostrofisch.

\textbf{4.} Tank $\num{e-6}/(1.257\times0.04) = \num{2.0e-5}$; depressie
$10/(\num{1.084e-4}\times10^8) = \num{9.2e-4}$; gyre
$\num{e-3}/(\num{1.084e-4}\times10^6) = \num{9.2e-6}$.

\textbf{5.} $\delta/H = \sqrt{2\nu/(fH^2)} = \sqrt{2\,\mathrm{Ek}}$. Tank
$\sqrt{2\times\num{e-6}/1.257} = \qty{1.3}{mm}$; lucht $\sqrt{20/\num{1.084e-4}}
= \qty{430}{m}$.

\textbf{6.} Wrijving werkt in lagen van enkele duizendsten van de diepte in de
tank en de oceaan, enkele honderdsten in de atmosfeer; daarboven verloopt
de stroming vrijwel wrijvingsloos.

\textbf{7.} $\rho = p/(RT) = 101\,290/(286.9\times288.1) =
\qty{1.225}{kg/m^3}$.

\textbf{8.} $u_g = (1/\rho f)\,|\partial p/\partial y| = (400/300\,000)/(1.225
\times\num{1.084e-4}) = \qty{10.0}{m/s}$, uit het westen, met de lage
druk in het noorden links.

\textbf{9.} $\mathrm{Ro}_g = 10.04/(\num{1.084e-4}\times600\,000) = 0.154$;
$V = 2\times10.04/(1 + \sqrt{1.618}) = \qty{8.84}{m/s}$, \qty{12}{\percent}
langzamer.

\textbf{10.} $V = 2\times10.04/(1 + \sqrt{0.382}) = \qty{12.4}{m/s}$. De
grootste $V_g$ is $fr/4 = \qty{16.3}{m/s}$, een gradiënt $\rho fV_g =
\qty{2.16e-3}{Pa/m}$, \qty{6.5}{hPa} per \qty{300}{km}.

\textbf{11.} $g/c_p = \qty{9.77}{K/km}$. $N^2 = (9.81/255.7)\times(9.77 -
6.49)\times10^{-3} = \qty{1.26e-4}{s^{-2}}$, $N = \qty{0.0112}{s^{-1}}$, een
periode van \qty{9.3}{min}.

\textbf{12.} Met $\Gamma = 0$, $N^2 = (g/T)(g/c_p) = g^2/(c_pT) =
\qty{4.43e-4}{s^{-2}}$, $N = \qty{0.0210}{s^{-1}}$, 1.9 keer de
troposferische waarde.

\textbf{13.} $L_R = 0.0112\times10^4/\num{1.084e-4} = \qty{1035}{km}$, de
grootte van de depressie.

\textbf{14.} $g' = \qty{0.0147}{m/s^2}$; $\sqrt{g'h} = \qty{0.54}{m/s}$;
$L_R = 0.542/\num{1.084e-4} = \qty{5.0}{km}$.

\textbf{15.} Ja: het meer, \qty{12}{km} breed, is meer dan twee keer zijn
\omterm{def:b3:gfd-foundations:rossby-radius}{Rossby-deformatiestraal}, dus interne golven en stromingen worden door de
corioliskracht.

\textbf{16.} $(249.4 - 253.0)/600 = \qty{-6.0}{K}$ per \qty{1000}{km}.

\textbf{17.} $\ln(1013/265) = 1.3409$. Zuid:
$286.9\times253.0\times1.3409/9.81 = \qty{9922}{m}$; noord \qty{9781}{m}.
Het oppervlak van \qty{265}{hPa} ligt \qty{141}{m} lager in het noorden, een helling van
\num{2.35e-4} ten opzichte van het onderste oppervlak.

\textbf{18.} $\Delta u = (g/f)\times\text{helling} = 9.81\times\num{2.35e-4}/
\num{1.084e-4} = \qty{21.3}{m/s}$.

\textbf{19.} $(R/f)\ln(p_b/p_t)\,|\partial\bar T/\partial y| =
(286.9/\num{1.084e-4})\times1.3409\times\num{6.0e-6} = \qty{21.3}{m/s}$, de
hetzelfde: de twee berekeningen zijn één.

\textbf{20.} \omterm{def:b3:gfd-foundations:baroclinic}{Baroklien}: de temperatuur, en dus de dichtheid, verandert langs
de drukvlakken. \omterm{def:b3:gfd-foundations:baroclinic}{Barotroop} zou de wind niet met de hoogte veranderen:
\qty{10.0}{m/s}.

\textbf{21.} $(g/f\bar T)\,|\partial T/\partial y|\,\Delta z = 9.81\times
\num{6.0e-6}\times10^4/(\num{1.084e-4}\times255.7) = \qty{21.2}{m/s}$,
binnen \qty{0.3}{\percent}.

\textbf{22.} $10.04 + (286.9/\num{1.084e-4})\times\num{6.0e-6}\times
\ln(1013/541) = 10.04 + 9.96 = \qty{20.0}{m/s}$.

\textbf{23.} $\partial u/\partial z = 21.3/10^4 = \qty{2.13e-3}{s^{-1}}$;
$\mathrm{Ri} = \num{1.26e-4}/\num{4.54e-6} = 28$, ruim boven 0.25: de windschuiving
maakt de vrije troposfeer niet turbulent.

\textbf{24.} De wind van Orrin, \qty{31}{m/s}, is sterker dan het zonaal gemiddelde
op \ang{45}~NB en zwakker dan op \ang{30}~NB. De straalstroom van het polaire front verschuift
noord- en zuidwaarts; op elk moment is de gradiënt ervan, en daarmee de wind,
geconcentreerd over een band, en middelen over een breedtecirkel smeert dat uit.

\textbf{25.} $10.0 + 21.3$: de geostrofische wind op \qty{10}{km} boven de
steppe is \textbf{\qty[text-series-to-math=true]{31}{m/s}}, uit het westen
(\qty{31.3}{m/s} volgens de berekening).
\end{solution}
